\documentclass[11pt]{article}
\usepackage{series}
\usepackage{amsmath,amssymb,amsthm}
\usepackage{geometry}
\usepackage{hyperref}
\usepackage{enumitem}

\geometry{margin=1in}

% ---------- metadata ----------
\title{The Modal Triplet Theory Program B0:\\ Why Description Forces Circle, Lens, and Nil\\
and Why the Minimal Continuous Realization Is Ten-Dimensional}
\author{Peter Nero}
\date{January, 2026}

\begin{document}

\maketitle

\begin{abstract}
We classify the possible ways in which global coherence can fail in any theory that admits
local reduced descriptions with overlap consistency and identity persistence.
Working purely at the structural level and without assuming geometry, bundles, dynamics,
or dimensionality, we show that there exist exactly three and only three obstruction types:
loop inconsistency (circle), redundancy of representation (lens), and failure of
representability (nil). These obstructions are exhaustive and mutually irreducible.

We further demonstrate that any realization capable of supporting all three obstruction
types simultaneously in a stable and non-degenerate manner must exhibit a tri-layer overlap
structure. When combined with continuity and kinematic persistence requirements, the
minimal such realization is equivalent, up to admissible re-encoding, to a ten-dimensional
configuration space composed of a four-dimensional effective base and three independent
internal layers. Geometry and bundle structure are shown to arise as the minimal bookkeeping
framework required to encode circle, lens, and nil consistently, rather than as fundamental
axioms.

This work provides the missing classification link between the structural core of Modal
Triplet Theory and its encoding-class and realization papers, explaining why gravity,
gauge redundancy, and termination of description are inevitable structural responses to
local describability and global obstruction.
\end{abstract}

\tableofcontents
\newpage

% ============================================================
\section{Introduction and Scope}

The structural core of Modal Triplet Theory establishes general conditions under which
reduced descriptions exist, how such descriptions relate on overlaps, and why no single
global reduced description can exist. These results are formulated abstractly, in terms of
projection, admissibility, and identity persistence, and make no assumptions about geometry,
spacetime, bundles, gauge structure, or dimensionality.

However, once local describability and global obstruction are taken seriously, a natural
question arises: \emph{what forms of failure of global coherence are actually possible, and
why do the same geometric and gauge-like structures appear so universally across effective
physical theories?} The purpose of the present paper is to answer this question at the level
of structural classification.

We show that the failure of global coherence in a locally describable system admits an
exhaustive and minimal classification. Any such failure must appear as one of three distinct
types: inconsistency detected only around closed overlap chains (circle), non-uniqueness of
local representation without contradiction (lens), or complete loss of representability
(nil). No fourth obstruction type is possible under the assumptions of local describability
and overlap consistency.

Having classified the obstruction types, we then address the question of realization.
We ask: \emph{what minimal kind of structure can support all three obstructions
simultaneously, without degenerating one into another?} We show that a tri-layer overlap
architecture is the minimal solution. Furthermore, when continuity and kinematic
persistence are required, this tri-layer structure admits a minimal continuous realization
equivalent, up to admissible re-encoding, to a ten-dimensional configuration space.

Throughout this work, geometry and bundle structure are treated not as fundamental inputs
but as emergent bookkeeping frameworks forced by the need to encode overlap consistency,
redundancy, and termination of description. The results of this paper therefore do not
postulate gravity, gauge symmetry, or dimensionality; instead, they explain why these
structures arise inevitably once local description and global obstruction are both present.

\begin{remark}[Relation to the series]
This paper depends only on the structural core of Modal Triplet Theory and its kinematic
interpretation of persistence across admissible overlaps. It precedes and motivates the
encoding-class papers on gravity, gauge structure, quantization, and unified encodings, as
well as the realization paper on geometric and bundle models. No realization-specific
assumptions are used in the classification arguments presented here.
\end{remark}

% ============================================================
% Next section will be:
% \section{Structural Assumptions and the Meaning of Global Coherence}
% ============================================================
\section{Structural Assumptions and Global Coherence}

In this section we state explicitly the structural assumptions imported from the core
of Modal Triplet Theory and clarify what is meant by global coherence and its failure.
No geometric, topological, or dynamical assumptions are introduced.

\subsection{Local describability}

\begin{assumption}[Local describability]
There exists a collection of admissible domains $\{A_\alpha\}$ covering the region of
interest such that on each $A_\alpha$ a reduced description is well-defined.
Equivalently, there exist local encodings
\[
\mathcal E_\alpha = (A_\alpha, Z_\alpha, E_\alpha, \ldots)
\]
with associated local sections selecting representatives of coherent structure.
\end{assumption}

Local describability asserts only that reduced descriptions exist \emph{locally};
it does not assume the existence of any globally valid reduced description.

\subsection{Overlap consistency}

\begin{assumption}[Overlap consistency]
On any nonempty overlap $A_{\alpha\beta} := A_\alpha \cap A_\beta$, the corresponding
local encodings are related by admissible re-encoding maps
\[
f_{\alpha\beta} : Z_\alpha \to Z_\beta
\]
that preserve coherent content up to admissible equivalence.
\end{assumption}

Overlap consistency ensures that different local descriptions of the same underlying
coherent structure do not contradict one another on shared domains.

\subsection{Identity persistence}

\begin{assumption}[Identity persistence]
If a coherent structure is representable on a chain of overlapping admissible domains,
then its identity can be tracked consistently across that chain, up to admissible
re-encoding equivalence.
\end{assumption}

This assumption underlies the kinematic notion of persistence developed in the
coherent kinematics paper. It does not assume trajectories, spacetime, or equations
of motion; it requires only that ``being the same'' is a meaningful notion locally.

\subsection{Finite admissibility and global obstruction}

\begin{assumption}[Finite admissibility]
There exists no single admissible domain that covers the entire region of interest.
Equivalently, no global section of the reduced description exists.
\end{assumption}

Finite admissibility is the structural statement that global coherence is generically
obstructed. It is not a statement about experimental limitation or lack of information,
but a property of the reduced description itself.

\subsection{Global coherence}

We now formalize the notion whose failure we seek to classify.

\begin{definition}[Global coherence]
A system is said to admit global coherence if there exists a single reduced encoding
$\mathcal E = (A, Z, E, \ldots)$ with $A$ equal to the entire region of interest, such
that all local descriptions arise as restrictions of $\mathcal E$.
\end{definition}

Global coherence therefore requires the existence of a global section compatible with
all local encodings and all overlap relations.

\begin{definition}[Failure of global coherence]
Failure of global coherence occurs when local describability and overlap consistency
hold, but no such global encoding exists.
\end{definition}

\subsection{The classification problem}

Given the assumptions above, the central question addressed in this paper is:

\begin{quote}
\emph{In a system that is locally describable, overlap-consistent, and admits identity
persistence but fails to be globally coherent, what distinct structural modes of failure
are possible?}
\end{quote}

The remainder of this work shows that this question has a sharp and exhaustive answer.
Under the stated assumptions, failure of global coherence must take one of exactly three
forms, which we classify in the following sections.

% ============================================================
% Next section will be:
% \section{Exhaustive Classification of Obstruction Types}
% ============================================================
% ============================================================
\section{Obstruction Types as Invariants of the Encoding Atlas}

In this section we formalize what is meant by an ``obstruction type'' and show
that failure of global coherence admits an exhaustive classification into exactly
three inequivalent classes. The classification is stated in terms of invariants
of the encoding atlas and is therefore stable under admissible refinement and
re-encoding.

\subsection{The encoding atlas and its nerve}

We briefly recall the structural objects introduced in the Modal Triplet Theory
core.

Let $\{A_\alpha\}$ be an admissible cover of the region of interest, with
associated local encodings
\[
\mathcal E_\alpha = (A_\alpha, Z_\alpha, E_\alpha, \ldots).
\]
On nonempty overlaps $A_{\alpha\beta}$ there exist admissible re-encoding maps
\[
f_{\alpha\beta} : Z_\alpha \to Z_\beta,
\]
defined up to admissible equivalence.

These data define a groupoid of encodings and a simplicial object, the
\emph{atlas nerve} $N$, whose $k$-simplices correspond to nonempty $(k+1)$-fold
overlaps of admissible domains.

\subsection{Global coherence and descent}

Global coherence is equivalent to the existence of a global object descending
from the local encodings.

\begin{definition}[Global object]
A \emph{global object} is a choice of representatives in each $Z_\alpha$
together with compatibility on all overlaps such that all re-encoding maps
compose consistently on the atlas nerve.
\end{definition}

Failure of global coherence is therefore a failure of descent: local data exist,
but cannot be glued into a global object.

\subsection{Obstruction types as atlas invariants}

We now define precisely what is meant by an obstruction \emph{type}.

\begin{definition}[Obstruction type]
An \emph{obstruction type} is an equivalence class of failures of global descent
that:
\begin{enumerate}[label=(\roman*)]
\item is invariant under admissible refinement of the cover;
\item is invariant under admissible re-encoding;
\item cannot be eliminated by local modification of representatives;
\item corresponds to a distinct failure mode of the descent problem.
\end{enumerate}
\end{definition}

Two failures of global coherence are of the same obstruction type if they define
the same invariant class on the atlas nerve or encoding fibration.

\subsection{Three invariant failure modes}

Under the assumptions of local describability, overlap consistency, identity
persistence, and finite admissibility, exactly three inequivalent obstruction
types can arise.

\begin{theorem}[Exhaustive obstruction classification]
Any obstruction to global coherence is equivalent, under admissible refinement
and re-encoding, to exactly one of the following:
\begin{enumerate}[label=(\roman*)]
\item \textbf{Circle:} a nontrivial holonomy class on a $1$-cycle of the atlas
nerve, corresponding to path-dependent inconsistency of transport;
\item \textbf{Lens:} a nontrivial isotropy or automorphism class of encoding
fibers, corresponding to non-unique but consistent local lifts;
\item \textbf{Nil:} a nontrivial complement of the encodable set
$X_{\mathrm{enc}} \subset X$, corresponding to empty encoding fibers.
\end{enumerate}
No fourth inequivalent obstruction type exists.
\end{theorem}

\subsection{Interpretation of the three types}

The three obstruction types correspond to distinct descent failures:

\begin{itemize}
\item Circle obstructs descent by path dependence on loops in the nerve.
\item Lens obstructs descent by non-uniqueness of lifts even when compatibility
holds.
\item Nil obstructs descent by failure of existence of local lifts.
\end{itemize}

These failure modes are mutually irreducible: circle cannot be absorbed into
lens, lens cannot be absorbed into nil, and nil cannot be resolved by refinement.

\begin{remark}
Formulated in this way, the classification is a theorem about the descent problem
for the encoding atlas. It does not rely on logical case-splitting, but on the
structure of the atlas nerve and the encoding fibration.
\end{remark}

% ============================================================
% The following sections specialize these three obstruction types:
% Section 4: Circle
% Section 5: Lens
% Section 6: Nil


\subsection{Classification theorem}

\begin{theorem}[Exhaustive obstruction classification]
Under the assumptions of local describability, overlap consistency,
identity persistence, and finite admissibility, any obstruction to global
coherence must be of exactly one of the following three types:
\begin{enumerate}[label=(\roman*)]
\item \textbf{Circle}: loop inconsistency detected only on closed chains of
overlapping admissible domains;
\item \textbf{Lens}: non-uniqueness of admissible local representation without
contradiction;
\item \textbf{Nil}: failure of representability, i.e.\ absence of any admissible
encoding on a region.
\end{enumerate}
No fourth obstruction type is possible.
\end{theorem}

\subsection{Proof sketch}

We sketch the argument, emphasizing structural necessity rather than
construction.

Consider an attempt to extend local descriptions to a global encoding.
Failure of this attempt must manifest in one of the following ways:

\paragraph{Case 1: Path-dependent inconsistency.}
Suppose that for every point there exists at least one admissible local
description, and that pairwise overlaps are consistent, but transporting a
coherent structure along two distinct chains of overlapping domains yields
inequivalent results. By identity persistence, such inconsistency can only be
detected when the two chains form a closed loop. This defines an obstruction
visible only on closed overlap chains, which we call a \emph{circle}.

\paragraph{Case 2: Redundant but consistent descriptions.}
Suppose that for every point there exist admissible local descriptions, and
that all overlap relations are consistent, but that the local description is
not unique. If multiple admissible sections represent the same coherent
structure with no canonical global choice, then global coherence fails by
non-uniqueness rather than contradiction. This defines a \emph{lens} obstruction.

\paragraph{Case 3: Failure of representability.}
Finally, suppose that there exist regions on which no admissible local
description applies. In this case, global coherence fails because the reduced
description ceases to exist. This defines a \emph{nil} obstruction.

These cases are mutually exclusive and collectively exhaustive. Any apparent
alternative failure mode reduces to one of them: path dependence reduces to
circle, non-canonical choice reduces to lens, and absence of description
reduces to nil. Therefore no fourth obstruction type exists.
\hfill$\square$

\subsection{Irreducibility of the three types}

It is important to note that the three obstruction types are irreducible to
one another under admissible re-encoding.

\begin{remark}
A circle obstruction cannot be removed by choosing a different local
representative, and therefore cannot be reduced to a lens. A lens obstruction
does not involve contradiction and therefore cannot be reduced to nil. Nil
represents the absence of description and therefore cannot be reduced to either
circle or lens. Each obstruction type represents a distinct structural failure
mode.
\end{remark}

\subsection{Interpretive note}

The terminology \emph{circle}, \emph{lens}, and \emph{nil} is descriptive rather
than metaphorical. Each name refers to the minimal structural pattern by which
the corresponding obstruction is detected: closed overlap chains, asymmetric
projection--reconstruction pairs, and empty encoding fibers, respectively.

% ============================================================
% Next section will be:
% \section{Circle: Loop Inconsistency and Holonomy}
% ============================================================
\section{Circle: Loop Inconsistency and Holonomy}

We now analyze the first obstruction type identified in the classification
theorem: \emph{circle}. This obstruction captures the failure of identity
persistence detected only on closed chains of overlapping admissible domains.

\subsection{Closed overlap chains}

Let $\{A_\alpha\}$ be an admissible cover, and consider a finite sequence of
domains
\[
A_{\alpha_1}, A_{\alpha_2}, \ldots, A_{\alpha_n}
\]
such that each consecutive overlap $A_{\alpha_i\alpha_{i+1}}$ is nonempty, and
$A_{\alpha_n\alpha_1}$ is also nonempty. Such a sequence defines a
\emph{closed overlap chain}.

On each overlap $A_{\alpha_i\alpha_{i+1}}$ there exists an admissible
re-encoding map
\[
f_{\alpha_i\alpha_{i+1}} : Z_{\alpha_i} \to Z_{\alpha_{i+1}}
\]
relating the local descriptions.

\subsection{Definition of circle}

\begin{definition}[Circle obstruction]
A \emph{circle} obstruction exists if there is a closed overlap chain such that
the composition of admissible re-encoding maps around the loop satisfies
\[
f_{\alpha_n\alpha_1} \circ f_{\alpha_{n-1}\alpha_n} \circ \cdots \circ
f_{\alpha_1\alpha_2}
\;\not\sim\;
\mathrm{id}_{Z_{\alpha_1}},
\]
where $\sim$ denotes admissible equivalence.
\end{definition}

Thus, although each pairwise overlap is consistent, identity persistence fails
when descriptions are transported around the loop.

\subsection{Minimality of triple overlaps}

The circle obstruction cannot be detected on pairwise overlaps alone.
At least three overlapping domains are required.

\begin{lemma}
If all admissible overlap chains have length at most two, then any apparent
loop inconsistency can be removed by admissible re-encoding, and no genuine
circle obstruction exists.
\end{lemma}

\begin{proof}[Sketch]
With only pairwise overlaps, any composition of re-encoding maps can be reduced
to a single transition between two charts. Any nontrivial effect can therefore
be absorbed into a redefinition of the local representative, reducing the
obstruction to a lens rather than a circle.
\end{proof}

This establishes that circle obstructions are intrinsically associated with
\emph{triple overlap structure}.

\subsection{Circle as shared obstruction}

A crucial feature of the circle obstruction is that it is not localized within
any single chart or encoding. Instead, it resides entirely in the compatibility
data between encodings.

\begin{remark}
The circle obstruction is \emph{shared}: it cannot be attributed to the base
description alone or to any internal layer alone. It arises from the failure of
simultaneous consistency of all overlap relations around a loop.
\end{remark}

This shared character is what later forces any bookkeeping structure to couple
base and internal descriptions universally.

\subsection{Circle and holonomy}

Although no geometric structure has been assumed, the formal pattern exhibited
by a circle obstruction is that of \emph{holonomy}: the outcome of transporting
identity around a loop depends on the path taken.

\begin{remark}
The appearance of holonomy-like behavior here is purely structural. It arises
from local describability and overlap consistency together with global
obstruction, and does not presuppose a connection, curvature, or manifold.
\end{remark}

In later papers, this structure will be shown to necessitate connection- and
curvature-like bookkeeping when one attempts to stabilize identity persistence
under refinement.

\subsection{Relation to other obstruction types}

It is important to distinguish circle from the other obstruction types.

\begin{itemize}
\item A circle obstruction involves contradiction detected only on closed
chains; it cannot be removed by choosing a different local representative.
\item A lens obstruction involves non-uniqueness without contradiction; no loop
is required.
\item A nil obstruction involves absence of any admissible description.
\end{itemize}

\begin{remark}
If a purported loop inconsistency can be eliminated by a redefinition of local
representatives, then the obstruction is not a circle but a lens. Genuine circle
obstructions are therefore precisely those that survive all admissible
re-encodings.
\end{remark}

\subsection{Preview of consequences}

The existence of circle obstructions has two immediate consequences that will
be developed in subsequent sections:
\begin{enumerate}
\item it forces any stable realization to include additional bookkeeping
structure to track loop-dependent inconsistency;
\item it couples base and internal descriptions, preventing their independent
global specification.
\end{enumerate}

These consequences underpin the later emergence of geometric and gravitational
encoding structures.

% ============================================================
% Next section will be:
% \section{Lens: Redundancy and Non-Unique Representation}
% ============================================================
\section{Lens: Redundancy and Non-Unique Representation}

We now analyze the second obstruction type identified in the classification
theorem: \emph{lens}. Unlike circle, which involves contradiction detected on
closed overlap chains, lens captures failure of global coherence through
non-uniqueness of local representation without inconsistency.

\subsection{Non-unique local sections}

Consider an admissible domain $A_\alpha$ with local encoding
\[
\mathcal E_\alpha = (A_\alpha, Z_\alpha, E_\alpha, \ldots).
\]
Local describability requires the existence of at least one admissible section
selecting a representative of coherent structure. However, admissibility does
not require this section to be unique.

\begin{definition}[Lens obstruction]
A \emph{lens} obstruction exists if there are multiple admissible local sections
or encodings representing the same coherent structure on a domain, such that:
\begin{enumerate}[label=(\roman*)]
\item all representations are mutually consistent on overlaps;
\item no canonical global choice of representative exists;
\item any two choices are related by admissible re-encoding.
\end{enumerate}
\end{definition}

Thus, lens represents redundancy of description rather than contradiction.

\subsection{Projection--reconstruction asymmetry}

The defining structural feature of lens is an asymmetry between projection and
reconstruction.

\begin{remark}
Lens obstructions arise whenever the reduced description involves a projection
that forgets information, together with a reconstruction (section) that is not
unique. Projection maps may be canonical, while reconstruction maps are
necessarily choice-dependent.
\end{remark}

This asymmetry ensures that multiple admissible representations coexist, all
equally valid, but none privileged.

\subsection{Lens as descriptive redundancy}

A lens obstruction does not indicate a failure of consistency. Instead, it
indicates that the same coherent structure admits multiple internal
representations.

\begin{remark}
Lens obstructions encode redundancy of description, not multiplicity of
underlying entities. They therefore represent a structural symmetry of the
encoding rather than a symmetry of the system being described.
\end{remark}

This distinction will later be essential for understanding gauge structure as an
encoding phenomenon rather than a fundamental physical symmetry.

\subsection{Minimality of lens}

Lens obstructions can occur already on pairwise overlaps and do not require
closed overlap chains.

\begin{lemma}
Lens obstructions can exist in the absence of circle obstructions.
\end{lemma}

\begin{proof}[Sketch]
Non-unique local sections may exist even when all overlap transition maps are
path-independent. In this case, redundancy persists but no loop inconsistency
arises.
\end{proof}

This distinguishes lens sharply from circle.

\subsection{Lens and irreducibility}

Lens obstructions are irreducible to the other obstruction types.

\begin{itemize}
\item A lens obstruction does not involve contradiction and therefore cannot be
reduced to a circle.
\item A lens obstruction does not involve absence of description and therefore
cannot be reduced to nil.
\end{itemize}

\begin{remark}
If redundancy can be removed by a single admissible global choice, then the
obstruction is not a genuine lens. A genuine lens exists precisely when
redundancy is unavoidable.
\end{remark}

\subsection{Preview of consequences}

The existence of lens obstructions has two key consequences that will be
developed later:
\begin{enumerate}
\item it necessitates bookkeeping of descriptive redundancy across overlaps;
\item it gives rise to gauge-like encoding structures when realizations are
introduced.
\end{enumerate}

Lens therefore captures the structural origin of gauge freedom without invoking
fields, connections, or group actions.

% ============================================================
% Next section will be:
% \section{Nil: Termination of Describability}
% ============================================================
\section{Nil: Termination of Describability}

We now analyze the third and final obstruction type identified in the
classification theorem: \emph{nil}. Unlike circle and lens, which arise from
inconsistencies or redundancies among admissible descriptions, nil represents
the complete failure of local describability itself.

\subsection{Absence of admissible encodings}

Nil arises when the conditions for admissibility fail entirely.

\begin{definition}[Nil obstruction]
A \emph{nil} obstruction exists on a region if there is no admissible local
encoding applicable to that region. Equivalently, the encoding fiber over that
region is empty.
\end{definition}

Nil therefore represents not inconsistency of description, nor multiplicity of
representation, but the absence of any reduced description.

\subsection{Nil versus non-existence}

It is essential to distinguish nil from ontological non-existence.

\begin{remark}
Nil does not assert that the underlying system ceases to exist. It asserts only
that no admissible reduced description applies. The obstruction is descriptive,
not ontological.
\end{remark}

This distinction is crucial for later interpretation of horizons, collapse, and
singular behavior.

\subsection{Structural origin of nil}

Nil obstructions arise whenever the admissibility conditions themselves are
violated, for example when:
\begin{itemize}
\item no local section exists;
\item overlap consistency fails irreparably;
\item contractivity or stability conditions break down;
\item truncation or projection errors exceed tolerance.
\end{itemize}

These failures cannot be repaired by re-encoding or refinement.

\subsection{Nil as a boundary object}

Nil plays the role of a boundary in the space of admissible descriptions.

\begin{remark}
Nil corresponds to the boundary of the encoding atlas: it is the locus at which
the atlas fails to admit any chart. In categorical terms, it acts as a zero
object for admissible encodings.
\end{remark}

This boundary character underlies the irreversibility associated with nil.

\subsection{Irreversibility and termination}

Once a nil obstruction is encountered, no continuation of the reduced
description is possible across that region.

\begin{lemma}
Nil obstructions are terminal: once encountered along a chain of admissible
descriptions, they cannot be passed through or undone by admissible re-encoding.
\end{lemma}

\begin{proof}[Sketch]
By definition, no admissible encoding exists in a nil region. Any attempt to
extend a reduced description across such a region would therefore violate local
describability.
\end{proof}

This terminal character distinguishes nil sharply from circle and lens.

\subsection{Nil and kinematic termination}

In the kinematic interpretation of persistence across admissible overlaps, nil
corresponds to termination of motion or history.

\begin{remark}
In coherent kinematics, nil appears as the endpoint of worldlines, horizons of
representability, or collapse of record persistence. These interpretations are
derived, not assumed.
\end{remark}

\subsection{Nil and irreducibility}

Nil obstructions are irreducible to the other obstruction types.

\begin{itemize}
\item Nil is not redundancy and therefore cannot be reduced to lens.
\item Nil is not contradiction and therefore cannot be reduced to circle.
\end{itemize}

\begin{remark}
If an apparent absence of description can be removed by refining the admissible
cover or choosing a different encoding, then the obstruction is not nil. A
genuine nil obstruction exists only when all admissible descriptions fail.
\end{remark}

\subsection{Preview of consequences}

Nil obstructions complete the classification of failure modes of global
coherence. Together with circle and lens, they will determine the minimal
architectural requirements for any realization capable of supporting local
describability, redundancy, and termination simultaneously.

% ============================================================
% Next section will be:
% \section{Minimal Coexistence: Why Three Layers Are Required}
% ============================================================
% ============================================================
\section{Minimal Coexistence: Why Three Layers Are Required}\label{sec:sec7}

Having classified the three obstruction types—circle, lens, and nil—as
inequivalent invariants of the encoding atlas, we now address the question of
\emph{coexistence}. Specifically, we ask what minimal architectural structure is
capable of supporting all three obstruction types simultaneously, in a stable
and non-degenerate manner.

\subsection{Structural meaning of a layer}

We begin by making precise what is meant by a \emph{layer} in this context.

\begin{definition}[Independent layer]
A \emph{layer} is an independent structural carrier of encoding data such that:
\begin{enumerate}[label=(\roman*)]
\item it admits local sections and overlap re-encodings;
\item its local automorphism group (lens redundancy) acts nontrivially on its
encoding data;
\item its transition data cannot be absorbed into the automorphism group of any
other single layer under admissible re-encoding.
\end{enumerate}
\end{definition}

The third condition expresses independence: loop or overlap defects associated
with one layer cannot be reinterpreted purely as redundancy in another.

\subsection{Single-layer degeneracy}

We first consider the possibility that all obstruction types could be supported
within a single layer.

\begin{lemma}
A single-layer architecture cannot support non-degenerate coexistence of circle,
lens, and nil obstructions.
\end{lemma}

\begin{proof}[Sketch]
In a single layer, any loop-dependent inconsistency detected on closed overlap
chains can always be absorbed into a redefinition of local representatives,
because the full transition structure acts within a single automorphism group.
Thus circle obstructions reduce to lens obstructions. Alternatively, failure of
representability collapses directly to nil. No independent circle obstruction
can persist.
\end{proof}

This shows that one layer is insufficient.

\subsection{Two-layer architectures and gauge-removability}

We next consider architectures with two independent layers.

\begin{lemma}
In a two-layer architecture satisfying the independence conditions above, any
loop obstruction is admissibly equivalent to a lens obstruction.
\end{lemma}

\begin{proof}[Sketch]
With two layers, transition data live in a product of two automorphism groups.
Any loop defect detected on a closed overlap chain can be decomposed into a pair
of layer-wise transformations. Because there are only two layers, the resulting
defect can always be absorbed into a redefinition of representatives in one
layer relative to the other. Equivalently, the loop holonomy lies entirely within
the combined isotropy group of the encoding fibers. Under admissible re-encoding,
such defects reduce to non-uniqueness of representation rather than genuine
path-dependent inconsistency.
\end{proof}

\begin{remark}
This result does not assert that two-layer systems cannot exhibit holonomy in an
abstract groupoid sense. It asserts that, under the MTT notion of admissible
re-encoding and layer independence, such holonomy is not an invariant obstruction
type and therefore does not constitute a genuine circle.
\end{remark}

Thus, two layers are insufficient to stabilize circle as an independent
obstruction.

\subsection{Tri-layer necessity}

We now show that three layers suffice and are minimal.

\begin{theorem}[Tri-layer minimality]
Three independent layers are the minimal structural architecture capable of
supporting non-degenerate coexistence of circle, lens, and nil obstructions.
\end{theorem}

\begin{proof}[Sketch]
With three independent layers, pairwise overlaps support lens obstructions via
layer-wise redundancy, while triple overlaps support genuine cocycle data that
cannot be absorbed into the automorphism group of any single layer. Loop defects
detected on three-layer overlap chains therefore survive admissible re-encoding
as invariant circle obstructions. Nil remains possible through simultaneous
failure of all layers. Minimality follows from the preceding lemmas.
\end{proof}

\subsection{Shared obstruction and universal coupling}

A key feature of the tri-layer architecture is that the circle obstruction is
\emph{shared} across all layers.

\begin{remark}
In a tri-layer system, the circle obstruction resides in the compatibility data
among all three layers simultaneously. It cannot be attributed to any single
layer and therefore couples all layers and the base description universally.
This shared character underlies the later emergence of geometric and
gravitational encoding structures.
\end{remark}

\subsection{Stability under refinement}

Finally, we note that the tri-layer architecture is stable under admissible
refinement.

\begin{remark}
Architectures with fewer layers fail to stabilize all three obstruction types
under refinement, while architectures with additional layers introduce
redundancy without generating new obstruction classes. Three layers therefore
represent the minimal stable architecture.
\end{remark}


% ============================================================
% Next section will be:
% \section{From Tri-Layer Structure to Minimal Continuous Realization}
% ============================================================
% ============================================================
\section{From Tri-Layer Structure to Minimal Continuous Realization}
\label{sec:sec8}

In this section we pass from the purely structural tri-layer result to a
rigorous statement about minimal \emph{continuous representability}.
To obtain a fully rigorous argument, we specify a realization class in which:
(i) encodings are modeled by smooth local trivializations,
(ii) overlap data defines Čech cocycles,
(iii) circle obstructions correspond to nontrivial holonomy (curvature) classes.

No claim of ontological dimensionality is made; the conclusion is a statement of
representability up to admissible re-encoding.

\subsection{Realization class and precise meaning of ``circle''}
\label{subsec:realization-class}

We fix a realization class sufficient to make the obstruction invariants
geometric (bundle-theoretic) objects.

\begin{assumption}[Smooth atlas realization]
\label{ass:smooth-atlas}
Let $A\subset X$ be an admissible domain. Assume that on $P(A)$ there exists a
finite or locally finite admissible cover $\{U_\alpha\}$ such that:
\begin{enumerate}[label=(\roman*)]
\item each $U_\alpha$ is a smooth manifold (or smooth submanifold of a common
ambient manifold);
\item each encoding chart on $U_\alpha$ is a smooth local trivialization with
transition maps $g_{\alpha\beta}$ on overlaps $U_{\alpha\beta}$;
\item the transition maps form a Čech 1-cocycle with values in a structure group
$G$ (possibly protocol-indexed), i.e.
\[
g_{\alpha\beta} g_{\beta\gamma} g_{\gamma\alpha} = e
\quad \text{on } U_{\alpha\beta\gamma},
\]
up to admissible equivalence.
\end{enumerate}
\end{assumption}

\begin{assumption}[Circle realized as nontrivial holonomy]
\label{ass:circle-holonomy}
Assume that, in this realization class, a circle obstruction corresponds to a
nontrivial holonomy/curvature class of a $G$-bundle with connection on $P(A)$.
Equivalently: there exists a principal $G$-bundle $P_G\to P(A)$ with a
connection whose curvature is not identically zero on $P(A)$.
\end{assumption}

\begin{remark}
Assumption~\ref{ass:circle-holonomy} is a realization hypothesis:
it specifies how the abstract circle invariant is represented in the smooth
realization class. It is the minimal standard way to make ``circle'' rigorous
as a stable invariant under refinement.
\end{remark}

\subsection{Why each independent layer requires at least two continuous coordinates}
\label{subsec:two-per-layer}

We now prove a lower bound on the continuous coordinate dimension required per
independent layer to support a stable circle obstruction in the realization
class above.

\begin{definition}[Layer carrier]
A \emph{layer carrier} is the minimal smooth manifold factor $B$ (in a local
product-type representation) on which one independent layer of transition data
acts nontrivially.
\end{definition}

\begin{theorem}[No stable circle on one-dimensional layer carriers]
\label{thm:no-circle-1d}
Under Assumptions~\ref{ass:smooth-atlas}--\ref{ass:circle-holonomy}, a
one-dimensional layer carrier cannot support a stable circle obstruction.
Consequently, any independent layer carrier must satisfy $\dim B \ge 2$.
\end{theorem}

\begin{proof}
Let $B$ be a connected smooth 1-manifold (possibly with boundary).
In the realization class, the circle obstruction is represented by a principal
$G$-bundle with connection whose curvature is nonzero somewhere on the domain.
Curvature is a horizontal $\mathfrak{g}$-valued 2-form $F\in \Omega^2(\cdot;\mathfrak g)$.
However, on a 1-manifold $B$ one has
\[
\Omega^2(B) = \{0\},
\]
since there are no nonzero differential 2-forms on a 1-dimensional manifold.
Therefore any connection restricted to a 1D carrier has identically vanishing
curvature. In particular, no nontrivial curvature/holonomy class can be supported
\emph{as a local, stable, infinitesimal invariant}. Any apparent loop effect on a
1D carrier is necessarily global and can be absorbed into representative choice,
hence reduces to a lens-type redundancy rather than a circle invariant in the
sense of Assumption~\ref{ass:circle-holonomy}.

Thus a one-dimensional carrier cannot realize the circle invariant as a stable
curvature class. Therefore $\dim B\ge 2$ is necessary for an independent layer.
\end{proof}

\begin{remark}
This argument is fully rigorous and uses only the standard fact that curvature is
a 2-form and 2-forms vanish on 1D manifolds. It does not assume any specific
dynamics.
\end{remark}

\subsection{Internal dimensional lower bound for tri-layer coexistence}
\label{subsec:internal6}

We now combine tri-layer minimality (Section~\ref{sec:sec7}) with
Theorem~\ref{thm:no-circle-1d}.

\begin{theorem}[Six-dimensional internal lower bound]
\label{thm:internal6}
Assume tri-layer minimality holds (Section~\ref{sec:sec7}) and that each layer
admits a smooth realization satisfying
Assumptions~\ref{ass:smooth-atlas}--\ref{ass:circle-holonomy}.
Then any continuous realization supporting non-degenerate coexistence of circle,
lens, and nil requires internal dimension at least six:
\[
\dim B_{\mathrm{internal}} \ge 3 \cdot 2 = 6.
\]
\end{theorem}

\begin{proof}
By tri-layer minimality, three independent layers are required.
By Theorem~\ref{thm:no-circle-1d}, each independent layer carrier has dimension
at least 2. Independence implies these carriers contribute additively to the
internal coordinate dimension. Hence
\[
\dim B_{\mathrm{internal}} \ge 2 + 2 + 2 = 6.
\]
\end{proof}

\subsection{Base dimension requirement as an explicit realization axiom}
\label{subsec:base4}

The minimal base dimension depends on what class of kinematic persistence one
requires. To keep this section fully rigorous without importing a separate
dimension-selection theorem, we state the base requirement as an explicit
realization assumption.

\begin{assumption}[Four-dimensional kinematic base]
\label{ass:base4}
In the realization class considered in this paper, kinematic persistence is
required in the sense of the coherent kinematics layer, and the effective base
manifold supporting that kinematics has dimension
\[
\dim Y = 4.
\]
\end{assumption}

\begin{remark}
Assumption~\ref{ass:base4} packages the kinematic dimension-selection content of
the kinematics program into the present paper. If a separate theorem-level
derivation of $\dim Y=4$ is supplied elsewhere, this assumption can be replaced
by a reference.
\end{remark}

\subsection{Minimal continuous representability: the ten-dimensional bound}
\label{subsec:tenD}

We now state and prove the dimensional representability bound.

\begin{theorem}[Minimal continuous representability by a ten-dimensional space]
\label{thm:tenD}
Assume:
\begin{enumerate}[label=(\roman*)]
\item tri-layer minimality (Section~\ref{sec:sec7});
\item the smooth atlas realization and circle-as-holonomy assumptions
(Assumptions~\ref{ass:smooth-atlas}--\ref{ass:circle-holonomy});
\item the four-dimensional kinematic base assumption (Assumption~\ref{ass:base4}).
\end{enumerate}
Then any continuous realization supporting non-degenerate coexistence of circle,
lens, and nil obstructions is \emph{representable}, up to admissible re-encoding,
by a configuration space of dimension at least ten:
\[
\dim M \;\ge\; \dim Y + \dim B_{\mathrm{internal}} \;\ge\; 4 + 6 \;=\; 10.
\]
Moreover, the bound is sharp in the realization class: it is achieved by a
product-type representation
\[
M \simeq Y_4 \times B_1 \times B_2 \times B_3,
\qquad \dim B_i = 2.
\]
\end{theorem}

\begin{proof}
By Assumption~\ref{ass:base4}, $\dim Y = 4$.
By Theorem~\ref{thm:internal6}, $\dim B_{\mathrm{internal}}\ge 6$.
Therefore any product-type representative has dimension at least $4+6=10$.

Sharpness: choose three independent layer carriers $B_i$ each of dimension 2.
Then $\dim(B_1\times B_2\times B_3)=6$ and hence $\dim M = 4+6=10$.
Such a representation is admissible in the realization class by construction,
and supports the required tri-layer obstruction structure.
\end{proof}

\begin{remark}[Representability, not ontology]
Theorem~\ref{thm:tenD} asserts representability up to admissible re-encoding
within the stated realization class. It does not assert that the physical world
is fundamentally ten-dimensional.
\end{remark}

\subsection{Bundle geometry as bookkeeping (consequence, not assumption)}
\label{subsec:bundle-bookkeeping}

Finally, we record the rigorous sense in which bundle geometry is selected as a
bookkeeping framework.

\begin{corollary}[Bundle geometry is the minimal bookkeeping framework]
\label{cor:bundle-bookkeeping}
Under Assumptions~\ref{ass:smooth-atlas}--\ref{ass:circle-holonomy}, the data of
local charts, overlap maps, and circle holonomy define a principal bundle with
connection on $P(A)$, unique up to admissible equivalence on the domain.
\end{corollary}

\begin{proof}
By Assumption~\ref{ass:smooth-atlas}, the transition maps form a cocycle defining
a principal $G$-bundle (standard Čech construction). By
Assumption~\ref{ass:circle-holonomy}, the circle invariant is represented by a
nontrivial connection curvature, so a connection exists on that bundle. Uniqueness
up to admissible equivalence follows from standard gauge equivalence of cocycles
and connections on overlaps.
\end{proof}

\begin{remark}
This corollary makes precise that bundle geometry is not postulated: it is the
canonical mathematical packaging of overlap consistency, redundancy, and holonomy
once a smooth realization class is chosen.
\end{remark}


% ============================================================
% Next section will be:
% \section{Consequences for Encoding Classes}
% ============================================================
\section{Consequences for Encoding Classes}

The results obtained in this work are structural and realizational, not
phenomenological. Their significance lies in constraining the possible
\emph{encoding classes} that can exist downstream in Modal Triplet Theory.
In this section we summarize how the circle--lens--nil classification and the
tri-layer minimality result determine the form of all subsequent encoding-class
papers.

\subsection{Encoding classes versus realizations}

An encoding class specifies a \emph{form of description}—a language in which
reduced structures may be expressed—while a realization specifies a concrete
implementation of that language.

\begin{remark}
The present paper constrains encoding classes but does not select particular
realizations. Bundle geometry, manifolds, connections, and spectral operators
appear only when encoding classes are realized concretely.
\end{remark}

The obstruction classification therefore precedes and constrains all encoding
papers, regardless of how they are later realized.

\subsection{Gravity as kinematic consistency encoding}

The circle obstruction represents loop-dependent inconsistency in identity
persistence. Any encoding class that stabilizes identity across closed overlap
chains must therefore introduce additional bookkeeping structure.

\begin{remark}
The encoding class commonly identified as gravity arises as the unique encoding
that resolves circle obstructions by enforcing consistency of kinematic
persistence across overlapping descriptions.
\end{remark}

This result does not assume geometry or field equations. It asserts only that
some encoding of consistency is unavoidable once circle obstructions exist.
Geometric and gravitational structures are realizations of this encoding class.

\subsection{Gauge structure as redundancy encoding}

Lens obstructions represent unavoidable redundancy of description.

\begin{remark}
Gauge structure arises as the encoding class that organizes and compensates for
lens obstructions. Gauge freedom is therefore redundancy of encoding, not a
symmetry of underlying reality.
\end{remark}

This explains the universality and unobservability of gauge degrees of freedom
across physical theories.

\subsection{Quantization as discrete constraint encoding}

Nil obstructions represent termination of describability.

\begin{remark}
Discrete spectra and quantization arise as encoding responses to nil
obstructions, in which continuous descriptions fail and only discrete,
topologically protected structures remain admissible.
\end{remark}

Quantization is therefore an encoding constraint, not an independent postulate.

\subsection{Unified and maximal encodings}

Encoding classes need not exist independently. In some regimes, multiple
obstruction types may be active simultaneously.

\begin{remark}
Encoding frameworks that simultaneously resolve circle, lens, and nil
obstructions represent maximal or saturated encodings. Such encodings naturally
exhibit features associated with unified frameworks, including coupled gravity,
gauge redundancy, and quantization constraints.
\end{remark}

The existence of such encodings is constrained by the tri-layer and minimal
dimensionality results derived above.

\subsection{Relation to dark sector interpretations}

The dark sector phenomena addressed elsewhere in the series are interpreted as
regimes in which only weaker encoding classes remain admissible.

\begin{itemize}
\item Statistical encoding classes dominate when local representability fails
but aggregate structure persists.
\item Relational encoding classes dominate when even statistical descriptions
fail and only transition constraints remain.
\end{itemize}

These interpretations rely directly on the obstruction classification and do not
introduce new structure.

\subsection{Position in the series}

This paper completes the structural classification required to motivate all
subsequent encoding-class papers in the Modal Triplet Theory program.

\begin{remark}
The logical sequence is as follows:
\begin{enumerate}[label=(\roman*)]
\item structural core establishes local describability and global obstruction;
\item coherent kinematics defines persistence and motion;
\item the present work classifies obstruction types and minimal coexistence;
\item encoding-class papers address gravity, gauge, quantization, and unification;
\item realization papers instantiate these encodings concretely.
\end{enumerate}
\end{remark}

\subsection{Outlook}

By showing that exactly three obstruction types exist and that their stable
coexistence forces a tri-layer architecture with a minimal continuous
realization, this work removes arbitrariness from the appearance of geometric,
gauge, and quantum structures in effective physical theories.

The remaining papers in the series build upon this classification to explore
specific encoding classes and their realizations, without revisiting the
structural inevitabilities established here.

% ============================================================
\section{Compatibility with the Modal Triplet Theory Core}

We briefly summarize how the obstruction classification developed in this work
maps directly onto the structural objects introduced in the core of Modal Triplet
Theory.

\subsection{Encoding atlas and obstruction types}

In the MTT core, admissible descriptions are organized into an encoding atlas,
with fibers
\[
\mathrm{Enc}(x)
\]
over points $x \in X$, and overlap relations encoded by admissible re-encodings.
Failure of global coherence corresponds to failure of descent in this atlas.

The three obstruction types classified here correspond to distinct invariant
features of this structure:

\begin{itemize}
\item \textbf{Circle} corresponds to a nontrivial holonomy class on $1$-cycles of
the atlas nerve $N$, detected as path dependence of admissible transport around
closed overlap chains.
\item \textbf{Lens} corresponds to nontrivial isotropy of encoding fibers
$\mathrm{Enc}(x)$, i.e.\ multiple admissible local sections related by
re-encoding automorphisms.
\item \textbf{Nil} corresponds to empty encoding fibers, defining the boundary of
the encodable region $X_{\mathrm{enc}} \subset X$.
\end{itemize}

These three invariants exhaust the possible failures of global descent in the
encoding atlas.

\subsection{Relation to the universal reduced relation}

The universal reduced relation
\[
\mathcal R \subset Y \times Y
\]
introduced in the MTT core provides a global object encoding admissible reduced
transitions. Circle obstructions correspond to nontrivial composition of
$\mathcal R$ along loops in the atlas nerve, lens obstructions correspond to
non-unique lifts of $\mathcal R$ to encoding fibers, and nil obstructions
correspond to regions where $\mathcal R$ admits no admissible lift at all.

\subsection{Role in the series}

The present paper refines the core statement that no global reduced encoding
exists by classifying \emph{how} such failure can occur. Subsequent papers in the
series treat gravity, gauge structure, and quantization as encoding responses to
circle, lens, and nil obstructions respectively, while realization papers
construct explicit geometric models satisfying the representability conditions
identified here.

No additional structural assumptions beyond those of the MTT core are introduced
in this classification.

% ============================================================
% End of main body
% ============================================================
\section{Conclusion}

In this work we have addressed a single structural question: given local
describability, overlap consistency, identity persistence, and the absence of a
global reduced description, what forms of failure of global coherence are
possible, and what minimal architectures can support them simultaneously?

We have shown that the answer is sharp. Exactly three and only three obstruction
types exist: loop inconsistency detected on closed overlap chains (circle),
redundancy of local representation without contradiction (lens), and termination
of describability (nil). These obstruction types are exhaustive and mutually
irreducible; no fourth mode of failure is compatible with the stated structural
assumptions.

We further demonstrated that non-degenerate coexistence of circle, lens, and nil
requires a tri-layer architectural structure. Architectures with fewer layers
collapse one obstruction type into another, while architectures with additional
layers introduce redundancy without increasing structural capacity. When
continuity and kinematic persistence are imposed as realization requirements,
this tri-layer structure admits a minimal continuous realization equivalent, up
to admissible re-encoding, to a ten-dimensional configuration space composed of a
four-dimensional effective base and three independent internal layers.

Throughout, geometry, bundle structure, and dimensionality have not been
assumed. Instead, they appear as minimal bookkeeping frameworks selected by the
need to encode shared loop obstructions, redundancy, and termination
consistently. The results therefore explain why geometric, gauge, and quantum
structures arise so universally in effective physical theories, without
postulating them as fundamental.

This paper completes the structural classification required to motivate the
encoding-class papers that follow in the Modal Triplet Theory series. Those works
develop gravity as a kinematic consistency encoding, gauge structure as a
redundancy encoding, quantization as a discrete constraint encoding, and unified
frameworks as saturated encodings, as well as concrete geometric and bundle
realizations. No additional structural assumptions are introduced downstream.

The classification presented here establishes that the familiar architectures of
modern theoretical physics are not arbitrary constructions, but inevitable
responses to the coexistence of local describability and global obstruction.

\end{document}
